
\subsection{Toward Lemma Z}\label{subsec:towardZ}
This subsection collects what is proved about Lemma~Z beyond
Theorem~\ref{thm:R0}.  Three recovery theorems are proved:
Theorem~\ref{thm:Bstar} (mates pinned on long windows are recovered at
every first row with finite base support),
Theorem~\ref{thm:Bpm} (signature room may be \emph{sign-mixed}: cofinite
contact sets are allowed) and Theorem~\ref{thm:S} (first rows whose
signature sets are \emph{exactly} swallowed by finitely many, or by
resonant and $d$-neutral, contact patterns are recovered without
approximating the first row).  The remaining cases are listed in
Problem~\ref{prob:lemmaZper} and Remark~\ref{rem:openZ}.  Several statements
are only sketched; they are labelled as such and are not used.

Local notation: $\phi_z$, the split fractions $\chi$, the sign-mixed rooms
$\mathfrak r_l,\mathfrak r^\star_l$, the products $\Lambda^\pm_f,
\Lambda^\star_f$, the swallowing signs $\epsilon_l$, the switching
amplitudes $\tau_l$, the weights $\mathfrak q_l$, the cost $\mathfrak c$
and the remainders $\mathfrak e_\pm,\mathfrak f^\pm_j,\mathfrak g_j$ are
local to this subsection.

\subsubsection*{Tools valid for every admissible operator}
In Lemmas~\ref{lem:phicalc}--\ref{lem:boundedfree} the operator $T$ is
admissible and otherwise arbitrary, $I$ is finite, $f\in S_{p^*}$ has
forced data as in Section~\ref{sec:setting}, $g\in\Cset(f)$, $t>0$, and
$(B_\pm,\Theta_\pm)$ is a two-sided decomposition of $g$ at scale $t$
(Lemma~\ref{lem:twosided}).  We write $\Delta B:=B_+-B_-$,
$\Delta\Theta:=\Theta_+-\Theta_-$ and
$\Delta\theta_{k,m}:=\lambda_{k,m}\Delta\Theta_m(k)$; since
$R_m^*W=\sum_k\lambda_{k,m}W(k)u_{k,m}$,
\begin{equation}\label{eq:DeltaB}
 \Delta B=-L^*\Delta\Theta=-\sum_{m\in I}\sum_k\Delta\theta_{k,m}u_{k,m},
\end{equation}
and for the operator of Definition~\ref{def:SLD}, $\Delta\theta_{k,m}=
\Delta\theta_{\mathfrak j(k,m)}$ in the notation of the previous subsections.
For $|z|\le1$ and $x\in\R$ put $\phi_z(x):=|x|-zx$.

\begin{lemma}[The functions $\phi_z$]\label{lem:phicalc}
Let $|z|\le1$, $x,y,r,c\in\R$ and $v\ge0$.
\begin{enumerate}
\item[(a)] $\phi_z(x)\ge(1-|z|)|x|\ge0$; if $|z|=1$, then $\phi_z(x)=2(zx)_-$
and $\phi_{-z}(x)=2(zx)_+$.
\item[(b)] $\phi_z(x-y)\le\phi_z(x)+\phi_{-z}(y)$ and
$|x|+|y|\le|x-y|+\phi_z(x)+\phi_{-z}(y)$.
\item[(c)] $|\phi_z(x+r)-\phi_z(x)|\le2|r|$.
\item[(d)] $\phi_z(-cv)=|c|v(1+z\sgn c)$.
\end{enumerate}
\end{lemma}

\begin{proof}
(a) $|x|-zx\ge|x|-|z||x|$; if $|z|=1$, $|x|=|zx|$.  (b) $\phi_z(x-y)=
|x-y|-zx+zy\le|x|-zx+|y|+zy$; and the difference of the two sides of the
second inequality is $|x-y|-z(x-y)\ge0$.  (c) $\big||x+r|-|x|\big|\le|r|$
and $|zr|\le|r|$.  (d) $|cv|+zcv=|c|v(1+z\sgn c)$.
\end{proof}

\begin{lemma}[Switching budget]\label{lem:switchbudget}
\[
 \sum_{j\notin F}\phi_{z_j}(B_+(j))\le\frac t{2q_0},\qquad
 \sum_{j\notin F}\phi_{-z_j}(B_-(j))\le\frac t{2q_0},\qquad
 \sum_{j\notin F}\phi_{z_j}(\Delta B(j))\le\frac t{q_0}.
\]
Consequently $\sum_{j\in K}2(z_j\Delta B(j))_-\le t/q_0$ and
$\sum_{j\notin F}(1-|z_j|)|\Delta B(j)|\le t/q_0$; in particular
$\sum_{j\in J_\gamma}|\Delta B(j)|\le t/(\gamma q_0)$ for $\gamma\in(0,1)$.
\end{lemma}

\begin{proof}
The derivation of \eqref{eq:budget} in the proof of
Lemma~\ref{lem:budget} uses only Lemma~\ref{lem:bookkeeping}(a),
$g(\xi)=0$ and $s(t)-1\le t^2/2$, so it holds for every $t>0$; as the
block excesses are nonnegative, $q_0G_b(a+tB_+)\le t^2/2$.  By
Lemma~\ref{lem:bookkeeping}(b), $G_b(a+tB_+)\ge\sum_{j\notin F}
(|tB_+(j)|-z_jtB_+(j))=t\sum_{j\notin F}\phi_{z_j}(B_+(j))$.  The $-$ side
is the same with $-B_-$ in place of $B_+$, and $\phi_z(-x)=\phi_{-z}(x)$.
The third inequality follows from Lemma~\ref{lem:phicalc}(b), the
consequences from Lemma~\ref{lem:phicalc}(a).
\end{proof}

\begin{remark}[What the budget does not say]\label{rem:budgetgloss}
The budget controls the $\ell_1$-mass of $\Delta B$ off $F$ only where
$z$ keeps room.  On contacts $\Delta B$ is $z$-signed up to $\ell_1$-mass
$t/(2q_0)$, but near contacts the $\ell_1$-mass of $\Delta B$ off $F\cup K$
need not be $O(t)$: if $z_j=1-1/j$ on an infinite set of coordinates
$j\notin F$, a switching $\Delta B(j)=c_j\ge0$ there is charged only
$\sum_jc_j/j$.  Only the $\phi_{z}$-weighted mass is $O(t)$ in general.
\end{remark}

\begin{lemma}[Split lemma]\label{lem:split}
Suppose $\Delta B1_{F^c}=V+e$ with $e\in\ell_1$ and $V\in\ell_1$ supported
in $K$ and $z$-signed \textup{(}$z_jV(j)\ge0$\textup{)}.  Then there are
$\chi:K\to[0,1]$ and $\mathfrak e_\pm\in\ell_1(F^c)$ with
\[
 B_+1_{F^c}=\chi V+\mathfrak e_+,\qquad B_-1_{F^c}=-(1-\chi)V+\mathfrak e_-,
 \qquad\|\mathfrak e_\pm\|_1\le\|e\|_1+\frac t{q_0}.
\]
\end{lemma}

\begin{proof}
Fix $j\notin F$ and put $x:=B_+(j)$, $y:=B_-(j)$, so $x-y=V(j)+e(j)$.  If
$V(j)=0$, let $\chi_j:=0$, $\mathfrak e_+(j):=x$, $\mathfrak e_-(j):=y$; by
Lemma~\ref{lem:phicalc}(b), $|x|,|y|\le|x|+|y|\le|e(j)|+\phi_{z_j}(x)+
\phi_{-z_j}(y)$.  If $V(j)\ne0$, then $|z_j|=1$; put $X:=z_jx$,
$Y:=z_jy$, $v:=z_jV(j)=|V(j)|>0$, so that $X-Y=v+z_je(j)$, and
$\phi_{z_j}(x)=2X_-$, $\phi_{-z_j}(y)=2Y_+$ by Lemma~\ref{lem:phicalc}(a).
Let $\chi_j:=\min(1,\max(0,X/v))$, $\mathfrak e_+(j):=x-\chi_jV(j)=
z_j(X-\chi_jv)$ and $\mathfrak e_-(j):=y+(1-\chi_j)V(j)=z_j(Y+(1-\chi_j)v)$.
Then $|\mathfrak e_+(j)|=\dist(X,[0,v])\le X_-+(X-v)_+\le X_-+Y_++|e(j)|$,
because $X-v=Y+z_je(j)$.  If $0<\chi_j<1$, then $|\mathfrak e_-(j)|=
|Y+v-X|=|e(j)|$; if $\chi_j=0$ ($X\le0$), then $|\mathfrak e_-(j)|=
|X-z_je(j)|\le X_-+|e(j)|$; if $\chi_j=1$ ($X\ge v$), then $|\mathfrak
e_-(j)|=|Y|$ and $Y=X-v-z_je(j)\ge-|e(j)|$, so $|Y|\le Y_++|e(j)|$.  In all
cases $|\mathfrak e_\pm(j)|\le|e(j)|+\phi_{z_j}(B_+(j))+\phi_{-z_j}(B_-(j))$;
sum over $j\notin F$ and use Lemma~\ref{lem:switchbudget}.
\end{proof}

Thus, when the switching difference is a $z$-signed contact vector up to
$e$, the $+$ side carries the fraction $\chi$ of it and the $-$ side the
complementary fraction with the opposite sign, coordinatewise.

\begin{lemma}[Peak pinning of the shift]\label{lem:peakshift}
Let $\eta\le\eta_*$, $0<t\le\min(t_\eta,1)$, and let $d_{\pm,m}$,
$\omega_{\pm,m}$ be as in Lemma~\ref{lem:suplevel};
put $\Delta d_m:=d_{+,m}-d_{-,m}$.  For every peak $k\in P_m$, with
$\varsigma_k:=\sgn w_m(k)$,
\begin{equation}\label{eq:peakshift}
 |\omega_{+,m}(k)|+|\omega_{-,m}(k)|=-\varsigma_k\Delta\Theta_m(k)-
 \Delta d_mM_m\ \ge0,
\end{equation}
and if $\alpha_m(k)\ne0$, the left side is at most
$t/(\sigma_m|\alpha_m(k)|)$, so that
\[
 -|\Delta\Theta_m(k)|-\frac t{\sigma_m|\alpha_m(k)|}\le\Delta d_mM_m\le
 |\Delta\Theta_m(k)|.
\]
Moreover
\begin{equation}\label{eq:didentity}
 \Delta d_mM_m=\frac1{mC_m}\sum_k\Phi_m(k)w_m(k)\Delta\theta_{k,m}+r_m,
 \qquad|r_m|\le\frac{2t}{\sigma_m}.
\end{equation}
\end{lemma}

\begin{proof}
Drop $m$.  By Lemma~\ref{lem:suplevel}(c), $\varsigma_k\omega_+(k)\le0\le
\varsigma_k\omega_-(k)$, so $|\omega_+(k)|+|\omega_-(k)|=\varsigma_k(\omega_-
-\omega_+)(k)$; and $\omega_--\omega_+=-\Delta\Theta-\Delta d\,w$ with
$w(k)=\varsigma_kM$.  This is \eqref{eq:peakshift}.  By
Lemma~\ref{lem:suplevel}(c) also $|\alpha(k)||\omega_\pm(k)|\le t/(2\sigma)$
for each sign, which gives the bound when $\alpha(k)\ne0$.  For
\eqref{eq:didentity}: by Lemma~\ref{lem:suplevel}(d), $\Delta d=d(\omega_+)-
d(\omega_-)+r'$ with $|r'|\le2t/\sigma$; since $d(w)=\|Dw\|_2^2/C=C$,
$d(\omega_+)-d(\omega_-)=d(\Delta\Theta)+C\Delta d$, hence $\Delta d\,M=
\ip{Dw}{D\Delta\Theta}/C+r'$, and $\ip{Dw}{D\Delta\Theta}=\sum_k\Phi(k)^2w(k)
\Delta\Theta(k)=\frac1m\sum_k\Phi(k)w(k)\Delta\theta_k$ because
$\lambda_k=m\Phi(k)$.
\end{proof}

Since $\|\alpha_m\|_1=1$, every block has a non-degenerate peak, so the
uniform shift $\Delta d_m$ is controlled by the switching at a single
non-degenerate peak; Proposition~\ref{prop:pinned}(c) used all
coordinates of the block.

\begin{lemma}[Flips on the base support]\label{lem:flip}
Here $F$ may be infinite.  For $j\in F$ put
\[
 \mathfrak f^+_j:=\Big(-\sgn(a_j)B_+(j)-\frac{|a_j|}t\Big)_+,\quad
 \mathfrak f^-_j:=\Big(\sgn(a_j)B_-(j)-\frac{|a_j|}t\Big)_+,\quad
 \mathfrak g_j:=\Big(\sgn(a_j)B_+(j)-\frac{2|a_j|}t\Big)_+ .
\]
Then $\sum_{j\in F}\mathfrak f^\pm_j\le t/(4q_0)$ and $\mathfrak g_j\le
|\Delta B(j)|+\mathfrak f^-_j$; and the clamp $b^{\rm cl}(j):=\sgn(B_+(j))
\min(|B_+(j)|,2|a_j|/t)$ $(j\in F)$ satisfies
\[
 \|(B_+-b^{\rm cl})1_F\|_1\le\|\Delta B1_F\|_1+\frac t{2q_0}.
\]
\end{lemma}

\begin{proof}
By Lemma~\ref{lem:bookkeeping}(b) and \eqref{eq:baseidentity}, the flip
part of $G_b(a+tB_+)$ is $\sum_{j\in F}2(-\sgn(a_j)tB_+(j)-|a_j|)_+=
2t\sum_j\mathfrak f^+_j$, and all other parts are nonnegative; as in the
proof of Lemma~\ref{lem:switchbudget}, $G_b(a+tB_+)\le t^2/(2q_0)$.  The
$-$ side ($-B_-$ in place of $B_+$) gives $2t\sum_j\mathfrak f^-_j\le
t^2/(2q_0)$.  If $\mathfrak g_j>0$, then $\sgn(a_j)B_+(j)=2|a_j|/t+\mathfrak
g_j$ and $\sgn(a_j)B_-(j)\le|a_j|/t+\mathfrak f^-_j$, so
$\sgn(a_j)\Delta B(j)\ge|a_j|/t+\mathfrak g_j-\mathfrak f^-_j\ge\mathfrak
g_j-\mathfrak f^-_j$.  Finally $|B_+(j)-b^{\rm cl}(j)|=(|B_+(j)|-2|a_j|/t)_+$,
which equals $\mathfrak g_j$ if $\sgn B_+(j)=\sgn a_j$ and is at most
$\mathfrak f^+_j$ otherwise; sum over $j\in F$.
\end{proof}

So near-flips on $F$ are pinned exactly like contacts: wherever $\|\Delta
B\|_1=O(Kt)$, clamping the base part on $F$ at $2|a_j|/t$ costs $O(Kt)$.

\begin{lemma}[Bounded free switching]\label{lem:boundedfree}
Let $F$ be finite, $\eta\le\eta_*$ with $\eta_\Gamma(\eta)\le1$ and
$0<t\le\min(t_\eta,1)$.  For every finite set $\mathcal U$ of pairs $(k,m)$,
$m\in I$, there is $C_{\mathcal U}<\infty$, depending only on $f$ and $\mathcal U$, such that
\[
 \max_{(k,m)\in\mathcal U}|\Delta\theta_{k,m}|\le C_{\mathcal U}\Big(1+\Big\|\sum_{(k,m)\notin\mathcal U}
 \Delta\theta_{k,m}u_{k,m}\Big\|_1\Big).
\]
\end{lemma}

\begin{proof}
Let $P^\perp$ be the orthogonal projection of $H$ onto $e^\perp$.  By
Lemma~\ref{lem:budget}(d), $q_0h(B_\pm)\le\Gamma_w(B_\pm,\Theta_\pm)\le2$,
i.e.\ $\|P^\perp U^*B_\pm\|\le c_f:=(2\nu/q_0)^{1/2}$.  By
\eqref{eq:DeltaB},
\[
 \Big\|P^\perp U^*\sum_{\mathcal U}\Delta\theta_{k,m}u_{k,m}\Big\|\le2c_f+\|U\|\,
 \Big\|\sum_{(k,m)\notin\mathcal U}\Delta\theta_{k,m}u_{k,m}\Big\|_1 .
\]
The linear map $c\mapsto P^\perp U^*\sum_{\mathcal U}c_{k,m}u_{k,m}$ on $\R^{\mathcal U}$ is
injective: if it vanishes, then $U^*\sum_{\mathcal U}c_{k,m}u_{k,m}$ is a multiple of
$e=U^*a/\nu$, so $\sum_{\mathcal U}c_{k,m}u_{k,m}=\mu a$ for some $\mu$ ($U^*$ is
injective); the left side lies in $Y$ and $a\in c_{00}$, so both sides
vanish by (T-c), and then $T\big(\sum_{\mathcal U}(c_{k,m}/q^*(Te_{k,m}))e_{k,m}
\big)=0$ forces $c=0$ by (T-b).  An injective linear map on a
finite-dimensional space is bounded below.
\end{proof}

For the operator of Definition~\ref{def:SLD} and a finite set $\mathcal U$ of ladder
indices, the lemma says: if all carriers outside $\mathcal U$ are pinned on a window
($\sum_{l\notin\mathcal U}|\Delta\theta_l|\le Kt$), then the switching through $\mathcal U$
is \emph{bounded}, not merely box-bounded ($|\Delta\theta_l|\le6\lambda_l/t$,
Lemma~\ref{lem:box}).  This is far from pinning, which on the window
$\mathcal W(l_*)$ requires $|\Delta\theta_l|\lesssim KT_{\rm lo}(l_*)$ at
the bottom scale.
